\scanpage{4}
\mspage{2}

\uline{Cor} Si $f$ est loc. de présentation finie et plat, $f$ est \inserted{\uline{univ}} \uline{ouvert}.

\uline{Prop} Si $S'\overset{f}{\longrightarrow}S$ est fid plat, l'application $T\longmapsto f^{-1}(T)$ des sous-schémas de $S$ vers les sous-schémas de $S'$ \deleted{est \unsure{injectif}} est injective.

\uline{\unsure{Cor|Prop}} Si $X,Y$ sur $S$, alors $\operatorname{Hom}_S(X,Y)\mathrel{-}\operatorname{Hom}_{S'}(X',Y')$ injective.

\deleted{Lemme}

car.

\diagram{\begin{tikzcd}[ampersand replacement=\&]
X \arrow[r,dash] \arrow[d,"f"'] \& X' \arrow[d,"f'"]\\
Y \& Y' \arrow[l,"\text{\uuline{plat}}"]
\end{tikzcd}}

\[
 f_*(F)'\overset{\sim}{\longrightarrow}f'_*(F')
\]

\diagram{\begin{tikzcd}[ampersand replacement=\&]
Z' \& Z \arrow[l] \\
S' \arrow[r,"f"'] \& S
\arrow[from=1-1,to=2-1,"g'"'] \arrow[from=1-2,to=2-2,"g\ \text{quasi-compact}"]
\end{tikzcd}}

$g(Z)=T$

\uline{Th} Soit $f:S'\longrightarrow S$ un morphisme plat. Pour tte partie \deleted{\unsure{la}} \deleted{\unsure{loc. fermée}} $T$ de $S$ \deleted{telle que} \inserted{image d'un morphisme} $Z\longrightarrow S$ quasi-compact, on a $f^{-1}(\overline T)=\overline{f^{-1}(T)}$.

\margin{\uline{Dém} \unsure{OPS|SPS} $X,Y,Z$ affines.}

\diagram{\begin{tikzcd}[ampersand replacement=\&]
C' \& C \arrow[l]\\
\unsure{I|J}\quad A' \arrow[u] \& A\quad I \arrow[l] \arrow[u,"u"]
\end{tikzcd}}

\margin{Alors $\overline T=V(I)$ ($I=\operatorname{Ker}u$) $f^{-1}(\overline T)=V(IA')$ \quad $\overline{f^{-1}(T)}=V(K)$ \quad $K=\operatorname{Ker}u$ \quad $\overline{g'(Z')}$}

\margin{\uline{Autre dém} On est ramené au \uline{cas} $g$ dominant $\Rightarrow g'$ dominant. \unsure{/ et|et} $\Leftarrow$ $g$ dominant, $S'\ne\varnothing\Longrightarrow Z'\ne\varnothing$. Mais cela est vrai pour \deleted{tt} morphismes \uline{générisants} car on est ramené aussitôt aux \uline{cor. Univ.} Soit $g:Z\longrightarrow S$ qu. cpt. Alors dominant \deleted{\unsure{sss|Mais}} pour tt pt maximal \deleted{\unsure{$s$}} $\unsure{\in S|\text{de }S}$ $Z_s\ne\varnothing$ \deleted{\unsure{est}}.}

\uline{Cor} Si $f$ fid plat \inserted{et quasi-compact}, la topologie de \deleted{$S$} \inserted{$f(S')$ induite par $S$} est quotient de celle de $S'$.

\deleted{$\unsure{T=f^{-1}(f(T))}\quad\overline T\quad f^{-1}(\overline{f(T)})\quad\unsure{\Longleftrightarrow}\quad\unsure{\text{si }}T=\overline T,\quad T\unsure{=}f^{-1}(f(T))$}

Il faut prouver que si $T'$ est une fermée de $S'$ saturée pour $f$, alors $T=f(Z')$ est rel. fermé dans $Z=f(S')$, i.e. que $\overline T\cap Z=T$ i.e. $\underset{\Vert\,T'}{f^{-1}(T)}=\underset{\Vert\,\overline{T'}}{f^{-1}(\overline T)}$ O.K.

\inserted{Si $f:S'\mathrel{-}S$ fid plat, qu. cpt \inserted{(ou mieux, gén., surjectif)}}

\uline{Cor}

\[
\begin{aligned}
 F(S)&\mathrel{-}F(S')= F(S'')\\
 O(S)&\mathrel{-}O(S')= O(S'')\\
 LF(S)&\mathrel{-}LF(S')= LF(S'')\\
 [LF^{\unsure{rc|qc}}(S)&\longrightarrow LF^{rc}= LF^{rc}(S'')\ ?\,]
\end{aligned}
\]

(où $F={}$ parties fermées \quad $O:{}$ parties ouvertes \quad $LF$ parties loc. fermées \inserted{rétrocompactes}) est \uline{exact} \quad $S''=S'\times_S S'$

% slanted
\margin{en fait aussi pour les parties \unsure{``constructibles''|structure induite} \deleted{\illegible{1}}}

\unsure{Bien Etd|Bien Ets} si $f$ surj, on a \quad $P(S)\mathrel{-}P(S')= P(S'')$ \quad ($P={}$ parties) car $S''\mathrel{-}|S'|\times_{|S|}|S'|$ surjectif. Donc l'exactitude des \deleted{deux} \inserted{deux} premiers exprime simplement que la top de $S$ est quotient, \deleted{\unsure{le}} \unsure{car|cas} est \unsure{bijectif|trivial|injective}.

\uline{Cor} Une partie \inserted{$T$} de $S$ est \unsure{ouverte|fermée} (fermée, loc fermée rétrocompacte) sss $f^{-1}(T)$ l'est. Si $T$ est l'image d'un morphisme qu. cpt, alors $T$ est loc. f. sss $f^{-1}(T)$ l'est\ldots

Pour le troisième cas, on est ramené aux deux premiers en introduisant l'adhérence et réutilisant le Th.

% archive sheet number: 3
